\documentclass{amsart}
% For dealing with references we use the comment environment
\usepackage{verbatim}
\newenvironment{reference}{\comment}{\endcomment}
%\newenvironment{reference}{}{}
\newenvironment{slogan}{\comment}{\endcomment}
\newenvironment{history}{\comment}{\endcomment}

% For commutative diagrams we use Xy-pic
\usepackage[all]{xy}

% We use 2cell for 2-commutative diagrams.
\xyoption{2cell}
\UseAllTwocells

% We use multicol for the list of chapters between chapters
\usepackage{multicol}

% This is generally recommended for better output
\usepackage{lmodern}
\usepackage[T1]{fontenc}

% For cross-file-references

% Package for hypertext links:
\usepackage{hyperref}

% For any local file, say "hello.tex" you want to link to please

% Theorem environments.
%
\theoremstyle{plain}
\newtheorem{theorem}[subsection]{Theorem}
\newtheorem{proposition}[subsection]{Proposition}
\newtheorem{lemma}[subsection]{Lemma}

\theoremstyle{definition}
\newtheorem{definition}[subsection]{Definition}
\newtheorem{example}[subsection]{Example}
\newtheorem{exercise}[subsection]{Exercise}
\newtheorem{situation}[subsection]{Situation}

\theoremstyle{remark}
\newtheorem{remark}[subsection]{Remark}
\newtheorem{remarks}[subsection]{Remarks}

\numberwithin{equation}{subsection}

% Macros
%
\def\lim{\mathop{\mathrm{lim}}\nolimits}
\def\colim{\mathop{\mathrm{colim}}\nolimits}
\def\Spec{\mathop{\mathrm{Spec}}}
\def\Hom{\mathop{\mathrm{Hom}}\nolimits}
\def\Ext{\mathop{\mathrm{Ext}}\nolimits}
\def\SheafHom{\mathop{\mathcal{H}\!\mathit{om}}\nolimits}
\def\SheafExt{\mathop{\mathcal{E}\!\mathit{xt}}\nolimits}
\def\Sch{\mathit{Sch}}
\def\Mor{\mathop{\mathrm{Mor}}\nolimits}
\def\Ob{\mathop{\mathrm{Ob}}\nolimits}
\def\Sh{\mathop{\mathit{Sh}}\nolimits}
\def\NL{\mathop{N\!L}\nolimits}
\def\CH{\mathop{\mathrm{CH}}\nolimits}
\def\proetale{{pro\text{-}\acute{e}tale}}
\def\etale{{\acute{e}tale}}
\def\QCoh{\mathit{QCoh}}
\def\Ker{\mathop{\mathrm{Ker}}}
\def\Im{\mathop{\mathrm{Im}}}
\def\Coker{\mathop{\mathrm{Coker}}}
\def\Coim{\mathop{\mathrm{Coim}}}

% Boxtimes
%
\DeclareMathSymbol{\boxtimes}{\mathbin}{AMSa}{"02}

%
% Macros for moduli stacks/spaces
%
\def\QCohstack{\mathcal{QC}\!\mathit{oh}}
\def\Cohstack{\mathcal{C}\!\mathit{oh}}
\def\Spacesstack{\mathcal{S}\!\mathit{paces}}
\def\Quotfunctor{\mathrm{Quot}}
\def\Hilbfunctor{\mathrm{Hilb}}
\def\Curvesstack{\mathcal{C}\!\mathit{urves}}
\def\Polarizedstack{\mathcal{P}\!\mathit{olarized}}
\def\Complexesstack{\mathcal{C}\!\mathit{omplexes}}
% \Pic is the operator that assigns to X its picard group, usage \Pic(X)
% \Picardstack_{X/B} denotes the Picard stack of X over B
% \Picardfunctor_{X/B} denotes the Picard functor of X over B
\def\Pic{\mathop{\mathrm{Pic}}\nolimits}
\def\Picardstack{\mathcal{P}\!\mathit{ic}}
\def\Picardfunctor{\mathrm{Pic}}
\def\Deformationcategory{\mathcal{D}\!\mathit{ef}}

\setlength{\emergencystretch}{3em}
\begin{document}
\title[Stacks Project, Tag 0AWY]{248 --- Gorenstein formal fibres of a local ring with a dualizing complex}
\maketitle
\noindent Copyright (C) 2005--2025 Johan de Jong. Stacks Project authors. Source: \href{https://stacks.math.columbia.edu/tag/0AWY}{Tag 0AWY}.

\noindent Editorial difficulty note (not author text). Difficulty H2: The argument reduces every formal fibre to the generic fibre of a domain quotient. It then tracks the dualizing complex through completion and fraction-field localization and proves the formal-fibre dualizing object is a shift of the unit. The compatibility of completion and the generic dualizing shift is a substantive structural obligation..

\noindent Editorial provenance note (not author text). History: excerpt from revision \texttt{a0444\allowbreak 6e57e\allowbreak c1fbc\allowbreak 252a8\allowbreak 71afc\allowbreak ec775\allowbreak 2fb28\allowbreak 07b14}, dualizing.tex, lines 4234--4264. Reference presentation was adapted; original-author context and core support blocks were retained or added. The mathematical wording is unchanged. Licensed under GNU FDL 1.2 or later; no invariant sections or cover texts. See ../licenses/COPYING. Context blocks reproduce author definitions and supporting results; they are not additional counted problems.

\section{Dualizing complexes}
\label{section-dualizing}

\noindent
In this section we define dualizing complexes for Noetherian rings.



\begin{lemma}
\label{lemma-dualizing-gorenstein-formal-fibres}
Let $A$ be a Noetherian local ring. If $A$ has a dualizing complex,
then the formal fibres of $A$ are Gorenstein.
\end{lemma}

\begin{proof}
Let $\mathfrak p$ be a prime of $A$. The formal fibre of $A$ at $\mathfrak p$
is isomorphic to the formal fibre of $A/\mathfrak p$ at $(0)$. The quotient
$A/\mathfrak p$ has a dualizing complex
(Lemma \href{https://stacks.math.columbia.edu/tag/0A7I}{lemma dualizing quotient (Tag 0A7I)}).
Thus it suffices to check the statement
when $A$ is a local domain and $\mathfrak p = (0)$.
Let $\omega_A^\bullet$ be a dualizing complex for $A$. Then
$\omega_A^\bullet \otimes_A A^\wedge$ is a dualizing complex
for the completion $A^\wedge$
(Lemma \href{https://stacks.math.columbia.edu/tag/0AWD}{lemma flat unramified (Tag 0AWD)}).
Then $\omega_A^\bullet \otimes_A K$ is a dualizing
complex for the fraction field $K$ of $A$
(Lemma \href{https://stacks.math.columbia.edu/tag/0A7G}{lemma dualizing localize (Tag 0A7G)}).
Hence $\omega_A^\bullet \otimes_A K$
is isomorphic to $K[n]$ for some $n \in \mathbf{Z}$.
Similarly, we conclude a dualizing complex for the formal fibre
$A^\wedge \otimes_A K$ is
$$
\omega_A^\bullet \otimes_A A^\wedge \otimes_{A^\wedge} (A^\wedge \otimes_A K) =
(\omega_A^\bullet \otimes_A K) \otimes_K (A^\wedge \otimes_A K) \cong
(A^\wedge \otimes_A K)[n]
$$
as desired.
\end{proof}
\end{document}
